\section{Grothendieck topologies}\label{SecGrothTop}
Let us start by recalling what a Grothendieck topology on an arbitrary category $\C$ is. First, recall that a \textit{sieve} on $\C$ is a full subcategory $S\subset \C$ such that, for every morphism $f : Y\to X$ in $C$, if $X$ belongs to $S$, then also $Y$ belongs to $S$. If $X\in C$, then \textit{a sieve on} $X$ is a sieve on the overcategory $C_{/X}$, which we usually denotes again by the letter $S$. More concretely, a sieve on $X$ is a right ideal of morphisms with codomain $X$, that is 
\begin{equation*}
	f\in S\implies f\circ g\in S, 
\end{equation*} 
for all $g$ such that the composition makes sense. 

Let $f: Y\to X$ be a morphism in the category $C$. If $S$ is a sieve on the object $X$, then we can define the \textit{pullback of the sieve} $S$ to be a sieve on $Y$ given by 
\begin{equation*}
	f^*(S):=\{g\in C_{/Y} \,| \, f\circ g\in S\}.
\end{equation*}
If $R=\{f_i:X_i\to X\}_{i\in I}$ is a collection of objects in $C_{/X}$, then there exists a smallest sieve $S\subset C_{/X}$ which contains each $f_i$. This is the full subcategory $S\subset C_{/X}$ given by
\begin{equation*}
	S:=\{f_ig \, |\, \mbox{dom}(f_i)=\mbox{cod}(g),\mbox{ for some }i\in I\}.
\end{equation*} 
We will say that $S$ is \textit{generated by the morphisms} $R=\{f_i\}_{i\in I}$, and we will write $S=(R)$. In particular, we will say that a sieve on $X$ is finitely generated if it is generated by a finite collection of morphisms in $C_{/X}$.
\begin{defi}\label{GTop}
	Let $C$ be a category. A \textit{Grothendieck topology} $\tau$ on $C$ is a map which assigns to each object $X\in C$ a \underline{collection} of sieves $\tau(C)$ such that, for each $X\in C$, the following conditions hold:
	\begin{itemize}[nosep]
		\item[(a)] the maximal sieve $C_{/X}$ is in $\tau(X)$;
		\item[(b)] if $S\in \tau(X)$ and $f:Y\to X$, then the pullback sieve $f^*(S)\in \tau(Y)$;
		\item[(c)] if $S\in \tau(X)$, and if $R$ is an arbitrary sieve on $X$ such that the pullback $f^*(R)\in \tau(Y)$, for every $f:Y\to X$ \underline{in $S$}, then $R\in \tau(X)$.
	\end{itemize} 
\end{defi}
In the following, if $S\in \tau(X)$, we will say that $S$ is a \textit{covering sieve} of $X$ or, simply, that $S$ \textit{covers} $X$. A category $C$ endowed with a Grothendieck topology will be called a \textit{site}, and will be denoted by $(C,\tau)$.
\begin{ex}\label{ExampleGrothOpenSub}
	The quintessential and motivating example of a Grothendieck topology is the one defined on topological spaces or, more precisely, on the category $\catname{Open}(X)$ of open subsets of $X\in \Top$. Indeed, a sieve $S$ over an open subset $U\in\catname{Open}(X)$ is nothing but a collection of open subsets of $U$ such that if $V'\subseteq V\in S$, then also $V'\in S$. A Grothendieck topology on $\catname{Open}(X)$ can be defined by requiring that a sieve covers the open subsets $U$ if $U$ is the union of all subsets in $S$.
\end{ex}
The above example shows an important detail. In general topology, an open cover of an open subset $U$ of $X$ is just a family $\{U_i \, |\, i\in I\}$ such that $U=\cup_i U_i$, but this is not necessarily a covering sieve as defined above. However, it does generate such a sieve for each $i\in I$, that is 
\begin{equation*}
	S:=\{V\subset U \, | \, V\subseteq U_i\}.
\end{equation*}
Noting that the pullback in $\catname{Open}(X)$ is just the intersection, we can generalize this situation.
\begin{defi}\label{GTopBasis}
	Let $C$ be a category with all pullbacks, a \textit{basis for a Grothendieck topology} is a map $k$ which assigns to each object $X\in C$ a \underline{collection} $k(X)$ of families of objects of	 $C_{/X}$ such that the following conditions hold:
	\begin{itemize}[nosep]
		\item[(a)] if $f:Y\to X$ is an isomorphism, then the singleton $\{f\}$ is in $k(X)$; 
		\item[(b)] if $\{f:Y_i\to X\}_{i\in I}\subset k(X)$, and $g:Y\to X$, then $\{Y\times_X Y_i\to Y\}_{i\in I}$ is in $k(Y)$;
		\item[(c)] if $\{f:Y_i\to X\ \, | \, i\in I\}\subset k(X)$, and for each $i\in I$ we have a family $\{g_{ij}:Z_j\to Y_i\, | \, j\in I_i\}$ in $k(Y_i)$, then the family $\{f_i\circ g_{ij}:Z_j\to X\, | \, i\in I, j\in I_i\}$ is in $k(X)$.
	\end{itemize} 
\end{defi}
In the following, an element of $k(X)$ will be called a \textit{covering family} of $X$. As the name suggests, if $k$ is a basis on $C$, then it generates a topology $\tau$ by
\begin{equation*}
	S\in \tau(X)\iff \exists R\in k(X)\mbox{ such that }R\subseteq S.
\end{equation*} 
Shortly, this means that a sieve is a covering sieve with respect to $\tau$ if and only if it contains a covering family of $k$ (as the case of topological spaces suggested). Suppose now we have a site $(C,\tau)$. First, note that, in general, there exist various basis for the topology $\tau$. However, among such we can find a maximal one, that is the basis $k$ given by
\begin{equation*}
	R\in k(X)\iff (R)\in \tau(X),
\end{equation*}
where $(R)$ is the sieve generated by the family $R$.
\begin{rem}
	In talking about Grothendieck topologies, it is customary (expecially in geometry) to consider basis rather then the actual topology. Therefore, when one defines a Grothendieck topology on a category $C$, she/he actually defines a basis for such topology. Furthermore, one often loosely says that a family $R=\{f_i:X_i\to X\}_{i\in I}$ covers $X$ if $(R)$ is a covering sieve in the topology. 
\end{rem}
\textcolor{red}{Mettere qualche esempio facile}

%==========================================================

\section{Sheaves for a Grothendieck Topology}\label{SheavesGrothTop}
As we have seen in Example \ref{ExampleGrothOpenSub}, Grothendieck topology represents a categorical extension of topology defined on a set. A useful construction in the latter setting is the one of sheaves on a topological space, so it is natural to ask how this notion can be extended to the case of a site. In what follows, we will consider a site $C=(C,\tau)$ and a presheaf $P:C^{op}\to \Set$. 
\begin{defi}
	Let $X\in C$ and let $S\in \tau(X)$ be a covering sieve for $X$, a \textit{matching family for $P$ with respect to $S$} \underline{is an element} of the limit
	\begin{equation*}
		\varprojlim_{Y\in S^{op}}P(Y).
	\end{equation*}
\end{defi}
\begin{rem}
	Some comments regarding the above definition is necessary. First of all, note that we adopted a customary abuse of notation in writing $Y\in S^{op}$, since we are actually considering the correspondent map $f\in S^{op}$ whose domain is $Y$. We considered the opposite of $S$ simply to correct for the fact that a presheaf $P$ is contravariant. Finally, writing the correspondent diagram it is easy to see that a matching family for $P$ with respect to $S$ is nothing but a collection of sections $\{y_f\}_{f\in S}$, each of which belongs to its correspondent $P(\mbox{dom}(f))$,  such that given another map $g:Z\to Y$ in $S$, it holds that
	\begin{equation*}
		P(g)(y_f)=y_{fg}.
	\end{equation*}
\end{rem}

For each object $X\in C$ and each covering sieve $S$ of $X$, there exists a canonical map 
\begin{equation}\label{canonicalMapSheaf}
	\theta_S:P(X)\to \varprojlim_{Y\in S^{op}}P(Y)	
\end{equation}
which restricts a global section to its subsections. Indeed, let $x\in P(X)$ be a section and define $\theta_S(x)=\{x_f\}_{f\in S}$ by restricting $x$ along each $f:Y\to X$, that is 
\begin{equation*}
	x_f:=P(f)(x)\in P(Y).
\end{equation*} 
This family is actually a matching family for $P$ with respect to $S$ because, given any composable morphism \( g : Z \to Y \) (so \( f \circ g \in S \)),
\[
P(g)(x_f)
= P(g)\big(P(f)(x)\big)
= P(f  g)(x)
= x_{f g}.
\]
Thus the family \( \{x_f\}_{f \in S} \) is indeed an element of the limit and we can therefore define
\[
\theta_S(x) := \big( P(f)(x) \big)_{f \in S}.
\]

\begin{defi}
	Let $X\in C$ and let $S\in \tau(X)$ be a covering sieve for $X$. If $m=\{x_f\}_{f\in S}$ is a matching family for $P$ with respect to $S$, an \textit{amalgamation} for $m$ is a section $x\in P(X)$ such that $x_f=P(f)(x)$, for each $f\in S$.
\end{defi}
The above discussion can be rephrased as follows. Given an object $X\in C$ and a covering sieve $S$ of $X$, every global section $x\in P(X)$ is an amalgamation for the matching family given by its subsections (specified by the covering sieve). However, recalling the local-to-global property of a sheaf over a topological space, what we would like to conclude is the inverse implication, \textit{i.e.} that given a matching family of $P$ with respect to $S$ we can always find an amalgamation, and we would like this amalgamation to be unique. All this boils down to the following definition. 
\begin{defi}
	A presheaf $P$ on a site $C$ is said to be a \textit{sheaf} if, for each object $X\in C$ and each covering sieve $S\in \tau(X)$, the canonical map \ref{canonicalMapSheaf} is a bijection. We will denote by $\catname{Sh}(C)$ the full subcategory of $\catname{Fun}(C^{op},\Set)$ consisting of sheaves over $C$.
\end{defi}
\begin{rem}
	The above definition of a sheaf over a site can be expressed in a form similar to the one usually given for a topological space. Recalling that in $\Set$ a limit is, by definition, an equalizer (see pag. 113 of \cite{maclane2013categories}) then the above discussion implies that $P$ is a sheaf on $C$ if, and only if, for each $X\in C$ and each covering sieve $S\in \tau(X)$, the diagram
	\begin{equation*}
		\begin{tikzcd}[column sep=small]% per separare le colonne
			P(X) \arrow[rr, "e"] &  & {\displaystyle\prod_{f\in S}P(\mbox{dom}(f))} \arrow[rr, "p", shift left] \arrow[rr, "q"', shift right] &  & {\displaystyle\prod_{f,g; f\in S}}P(\mbox{dom}(g))
		\end{tikzcd}
	\end{equation*}  
	is an equalizer of sets. Here $e(x)=\{P(f)(x)\}_{f\in S}$, the map $f$ ranges over $S$ and $g$ are maps composable with $f$ (hence $fg\in S$), while $p(x)_{f,g}=P(fg)(x)$ and $q(x)_{f,g}=P(g)(x_f)$.
\end{rem}
\begin{ex}
	If we consider the site consisting on the category $\catname{Open}(X)$, for a given topological space $X$, together with the Grothendieck topology of Example \ref{ExampleGrothOpenSub}, we obtain the category of sheaves on $X$, denoted as $\catname{Sh}(X)$.
\end{ex}
If the categoty $C$ is small, we can extend the following result, which holds for sheves on a topological space $X$.
\begin{thm}
	Let $C$ be a small category endowed with a Grothendieck topology $\tau$. Then, there exists an adjunction
	\begin{equation}\label{AdjSheaves}
		L:\catname{Fun}(C^{op},\Set)\leftrightarrows \catname{Sh}(C):\iota,
	\end{equation}
	where $\iota: \catname{Sh}(C)\to\catname{Fun}(C^{op},\Set)$ is the inclusion, while $L:\catname{Fun}(C^{op},\Set)\to\catname{Sh}(C)$ is the \textit{sheafification functor}. Moreover, $L$ preserves finite limits.
\end{thm}
\begin{proof}
	See \cite[Par. 3.5]{maclane2012sheaves}
\end{proof}
Recall that if $C$ is a locally small category, the Yoneda embedding $C\ni X\mapsto h_X(\_)=\Hom_C(\_,X)\in \catname{Fun}(C^{op},X)$ allows us to see the category $C$ as a particular full subcategory of its category of presheaves. Given a Grothendieck topology on $C$, we could ask if the correspondent presheaf $h_X$ are actually a sheaf, for each $x\in C$. This situation deserves a name. 
\begin{defi}
	A Grothendieck topology on a locally small category $C$ is said to be \textit{subcanonical} if, for every objects $X\in C$, the functor $h_X(\_)=\Hom_C(\_,X)$ is a sheaf.
\end{defi}
Note that if $C$ is a locally small category endowed with a subcanonical Grothendieck topology, then the Yoneda embedding is actually a fully faithful embedding $h:C\to \catname{Sh}(C)$.
\begin{rem}
	If $C$ is a locally small category endowed with a non-subcanonical Grothendieck topology, by applying the functor $L$ of Theorem \ref{AdjSheaves} to the presheaf $h_X$, for $X\in C$, we can define the correspondent sheaf, which we denote by $\tilde{h}_X$. Thus, we can also define a \textit{sheafified Yoneda embedding} by $X\mapsto\tilde{h}_X$. However, note that this functor is not fully faithful, it is only if the Grothendieck topology is subcanonical.
\end{rem}
We now discuss three important examples of Grothendieck topologies and sheaves which, in the meantime, allows us to introduce notions that will be used extensively in the following section.

%==========================================================

\subsection{The Regular Topology}
\begin{defi}
	Let $C$ be a category which admits pullbacks. We say that a morphism $f:X\to Y$ in $C$ is an \textit{effective epimorphism} if it exhibits $Y$ as the coequalizer of the projection maps $\pi_1,\pi_2:X\times_Y X\rightrightarrows X$.
\end{defi}
\begin{rem}
	Clearly, every effective epimorphism is also an epimorphism, while the converse need not be true. In particular, note that while $f:X\to Y$ being an epimorphism means that the induced map $\Hom_C(Y,Z)\to \Hom_C(X,Z)$ is injective for each $Z\in C$ (so that $\Hom_C(Y,Z)$ is actually a subclass of $\Hom_C(X,Z)$), being an effective epimorphism allows us to characterize this subclass, \textit{i.e.} $\Hom_C(Y,Z)$ are given by those maps $g:X\to Z$ such that the following diagram commutes
	\begin{equation*}
		\begin{tikzcd}
			& X \arrow[rd] \arrow[rrd, "g"]  &                     &   \\
			X\times_Y X \arrow[rd] \arrow[ru] &                                & Y \arrow[r, dashed] & Z \\
			& X \arrow[ru] \arrow[rru, "g"'] &                     &  
		\end{tikzcd}
	\end{equation*}
\end{rem}
\begin{lem}\label{Lemma:effectiveEpiComp}
	Let $C$ be a category which admits pullbacks, and suppose effective epimorphisms in $C$ are stable by pullbacks. Then, effective epimorphisms in $C$ are stable by compositions.
\end{lem}
\begin{proof}
	da scrivere
\end{proof}
\begin{defi}
	A category $C$ which admits pullbacks is said to be \textit{regular} if 
	\begin{itemize}[nosep]
		\item[(a)] it admits finite limits;
		\item[(b)] every morphism $f:X\to Y$ can be decomposed as $f=m\circ e$, where $e$ is an effective epimorphism and $m$ a monomorphism;
		\item[(c)] the collection of effective epimorphisms is stable under pullbacks.
	\end{itemize}
\end{defi}

\begin{defi}
	Let $C$ be a regular category and let $X\in C$. A sieve $S\in C_{/X}$ is a \textit{regular sieve} if it contains an effective epimorphism.
\end{defi}
Let $C$ be a regular category. Exploiting the above results and noting that, for each object $X$ in $C$, the identity morphism $\mbox{id}_X$ is an effective epimorphism, we can easily conclude the following result
\begin{lem}
	If $C$ is a regular category, the collection of regular sieves determines a Grothendieck topology on $C$. We will call this topology the \textit{regular topology} on $C$.
\end{lem} 

%==========================================================

\subsection{The Extensive Topology}
\begin{defi}
	Let $C$ be a category which admits pullbacks, and let $X,Y\in C$ such that $X\amalg Y$ exists in $C$. We say that $X\amalg Y$ is a \textit{disjoint coproduct} of $X$ and $Y$ if it holds that:
	\begin{itemize}[nosep]
		\item[(a)] both the inclusion maps $\iota_X:X\to X\amalg Y$ and $\iota_Y:Y\to X\amalg Y$ are monomorphisms;   
		\item[(b)] the pullback $X\times_{X\amalg Y}Y$ is an initial object of $C$.
	\end{itemize}
\end{defi}
\begin{ex}
	mettere un esempio/osservazione.
\end{ex}
\begin{defi}
	A category $C$ which admits finite limits is said to be \textit{extensive} if 
	\begin{itemize}[nosep]
		\item[(a)] it admits finite coproducts and coproducts in $C$ are disjoint;
		\item[(b)] for every morphism $f:X\to Y$ in $C$, the pullback functor
		\begin{equation*}
			f^*:C_{/Y}\to C_{/X},
		\end{equation*} 
		which sends a map $\alpha:Z\to Y$ in $f^*(\alpha):Z\times_Y X\to X$, preserves finite coproducts.
	\end{itemize}
\end{defi}

%==========================================================

\subsection{The Coherent Topology}

%==========================================================

\section{Grothendieck topoi and geometric morphisms}
\begin{defi}
	A category $\C$ is said to be a \textit{Grothendieck topos} if there exists a site $(C,\tau)$ such that $\mathcal{C}$ is equivalent to $\catname{Sh}(C)$. 
\end{defi}
The above definition is of limited practical use. However, a remarkable theorem of J. Giraud shows that Grothendieck topoi may be completely characterized by categorical requirements on $\C$. In order to state this result, we need some preliminary definitions.